\chapter{Especificação de Requisitos} \label{cap:especificacao}

Os requisitos abaixo descrevem a ferramenta como um todo, da entrada em código
C ao relatório ordenado que o revisor lê. Eles derivam do objetivo enunciado na
Seção~\ref{sec:objetivos} e das limitações observadas nas ferramentas
existentes, e alguns deles são restrições deliberadas, que proíbem
comportamentos que melhorariam os números finais ao custo do recall.

\section{Princípio que governa a especificação}
\label{sec:principio}

Uma única regra atravessa todos os requisitos. Um solucionador pode descartar
um candidato; o modelo de linguagem, não. Um veredicto de insatisfazibilidade é
uma prova de que a intercalação não existe, e o julgamento de um \llm\ não é
prova alguma. A consequência é que as etapas automáticas de filtragem
podem reduzir o conjunto, enquanto a etapa de triagem apenas o ordena e o
explica, mantendo tudo no relatório final.

A regra não decorre apenas de prudência. A literatura mede o efeito das duas
políticas no mesmo ponto de integração, e os filtros que descartam candidatos
perdem de 6 a 25 pontos de recall \cite{reducing-fp,adataint}, ao passo que a
política conservadora do LLift alcança recall de 1,00 \cite{llift}. Dela decorre
também a escolha das métricas, uma vez que a precisão deixa de ser a medida
principal e cede lugar ao recall, que precisa ser de 100\%, acompanhado do
\textit{Inspection Ratio} (Seção~\ref{sec:soundness}).

\section{Requisitos funcionais}
\label{sec:rf}

\begin{quadro}[htb]
  \caption{\label{quadro:rf}Requisitos funcionais}
  \small
  \begin{tabular}{|l|L{11.6cm}|}
    \hline
    & Descrição \\
    \hline
    RF1 & Aceitar como entrada um programa em C, com um ou vários arquivos, e um arquivo de configuração que declare os fluxos, suas prioridades e as primitivas de mascaramento da plataforma. \\
    \hline
    RF2 & Identificar as posições de memória alcançadas por mais de um fluxo, incluindo variáveis locais compartilhadas por meio de ponteiros. \\
    \hline
    RF3 & Enumerar todos os acessos de leitura e escrita a essas posições, registrando fluxo, função, caminho de chamadas, linha do código-fonte e laços que envolvem o acesso. \\
    \hline
    RF4 & Derivar candidatos a defeito a partir dessa enumeração, restringindo as triplas aos quatro padrões não serializáveis do Quadro~\ref{quadro:padroes}. \\
    \hline
    RF5 & Descartar candidatos cujos fluxos comprovadamente não podem se sobrepor, por prioridade ou por mascaramento. \\
    \hline
    RF6 & Verificar a viabilidade de caminho dos candidatos restantes com um solucionador SMT, descartando apenas os provados impossíveis. \\
    \hline
    RF7 & Emitir, para cada candidato sobrevivente, um registro de contexto com as evidências necessárias para julgá-lo. \\
    \hline
    RF8 & Responder a pedidos de contexto adicionais feitos durante a triagem, como a definição de uma função, o corpo de uma \isr\ ou o estado de mascaramento numa linha. \\
    \hline
    RF9 & Classificar cada candidato em um de quatro grupos de leitura, com explicação, sem remover nenhum candidato do relatório. \\
    \hline
    RF10 & Propor, para cada candidato, uma correção no vocabulário de interrupções, acompanhada de uma testemunha da intercalação que a justifica. \\
    \hline
    RF11 & Reverificar o programa corrigido, confirmando o desaparecimento do defeito e a ausência de defeitos novos. \\
    \hline
    RF12 & Oferecer um painel que permita configurar o programa, executar a análise, ler cada candidato ao lado do código-fonte e aceitar ou rejeitar uma correção. \\
    \hline
  \end{tabular}
  \fonte{Elaborado pelos autores.}
\end{quadro}

Três requisitos merecem comentário, porque neles está a diferença em relação às
ferramentas existentes.

O RF5 e o RF6 são requisitos de descarte, e ambos são condicionados a uma
prova. No RF5, a prova é estrutural, pois um fluxo de prioridade estritamente
menor não interrompe outro, e uma interrupção mascarada em todos os caminhos
não ocorre. No RF6, a prova é o veredicto de insatisfazibilidade do
solucionador. Um resultado inconclusivo, por limite de tempo ou por construção
não suportada, nunca é tratado como descarte, e o candidato segue adiante
carregando a informação de até onde o solucionador chegou.

O RF9 fixa o papel do modelo de linguagem, que ordena e explica sem alterar a
composição do relatório. Nisso reside a diferença entre reduzir o esforço de
revisão e reduzir o recall, e é também o que permite à avaliação adotar o
\textit{Inspection Ratio} como medida de custo.

O RF12, comumente tratado como acessório em trabalhos acadêmicos, é aqui parte
da contribuição. A comparação do Quadro~\ref{quadro:capacidades} mostra que a
automação é o eixo em que as ferramentas dessa literatura efetivamente
competem, com o NIChecker exigindo que o usuário indique variável e padrão a
cada execução e o IntRace exigindo um arquivo de configuração escrito
manualmente. Uma ferramenta que produz bons candidatos, mas não pode ser
operada, não atende ao objetivo a que se propõe.

\section{Requisitos não funcionais}
\label{sec:rnf}

\begin{quadro}[htb]
  \caption{\label{quadro:rnf}Requisitos não funcionais}
  \small
  \begin{tabular}{|l|L{11.6cm}|}
    \hline
    & Descrição \\
    \hline
    RNF1 & Nenhum defeito anotado no gabarito pode deixar de aparecer no relatório, e nenhum pode ser classificado no grupo de menor prioridade de leitura. \\
    \hline
    RNF2 & As premissas sob as quais o RNF1 vale devem estar declaradas por escrito e versionadas junto com o código. \\
    \hline
    RNF3 & Toda decisão de descarte deve ser rastreável até a prova que a autorizou. \\
    \hline
    RNF4 & O registro de contexto deve distinguir o que foi provado, o que veio da configuração e o que é desconhecido. \\
    \hline
    RNF5 & A etapa de triagem deve ser independente do provedor do modelo, permitindo trocar de modelo sem alterar os \textit{prompts}. \\
    \hline
    RNF6 & Toda medição deve registrar o modelo, a data e a configuração de \textit{prompt} que a produziram, e as respostas do modelo devem ser guardadas em cache para que uma medição possa ser recalculada sem nova chamada. \\
    \hline
    RNF7 & O código analisado deve entrar no \textit{prompt} delimitado e declarado como dado, nunca como instrução. \\
    \hline
    RNF8 & O custo por candidato, em \textit{tokens} e em tempo, deve ser medido e reportado. \\
    \hline
  \end{tabular}
  \fonte{Elaborado pelos autores.}
\end{quadro}

O RNF7 responde a um risco apontado por \citeonline{sastgenius}. O código
analisado entra no \textit{prompt} acompanhado de seus comentários, e um
comentário pode ter sido escrito para instruir o modelo a ignorar um defeito.
Num analisador de segurança, trata-se de um vetor de ataque; neste trabalho,
é ao menos uma via pela qual o modelo pode ser induzido ao erro por um
comentário desatualizado.

O RNF5 e o RNF6 existem porque comparações entre modelos envelhecem
rapidamente. O que este trabalho pretende afirmar não é que um modelo
específico funciona, e sim que uma determinada arquitetura de \textit{prompt}
funciona, afirmação que só é verificável se o modelo puder ser substituído com
o restante mantido fixo.

\section{Premissas e ameaças ao recall}
\label{sec:premissas}

O RNF1 é uma afirmação condicional, e as condições são as premissas abaixo.
Elas valem tanto como especificação quanto como declaração de limites.

\begin{itemize}
  \item O arquivo de configuração nomeia todos os fluxos. Um tratador não
    declarado é invisível para a análise, e esta é a maior ameaça ao recall.
  \item Uma chamada indireta cujos alvos não se resolvem pode chamar qualquer
    função de assinatura compatível cujo endereço tenha sido tomado.
  \item Posições compartilhadas são identificadas por análise \textit{may}, e a
    correção da análise de ponteiros usada é herdada, não reverificada.
  \item Uma função sem corpo disponível pode ler e escrever qualquer variável
    global.
  \item Fluxos de mesma prioridade podem se interromper mutuamente, e o
    aninhamento é permitido.
  \item Uma interrupção só está mascarada num intervalo se estiver mascarada em
    todos os caminhos e se nenhum outro fluxo puder reabilitá-la ali dentro.
  \item Apenas uma prova descarta um candidato, de modo que nem o resultado
    inconclusivo do solucionador nem o julgamento do \llm\ o fazem.
\end{itemize}

Uma premissa afasta-se deliberadamente desse padrão conservador, a de que uma
\isr\ não interrompe a si mesma. O Racebench não exclui a reentrância e a
literatura diverge quanto a ela, de modo que a escolha fica registrada como
perda potencial de recall, a ser revista antes de qualquer afirmação final.
